A central assumption in economics is that individuals make decisions with full information. In practice, however, acquiring and processing information requires time, attention, and cognitive effort. These informational frictions are particularly relevant for household electricity consumption, where the benefits of use are immediate but the costs are delayed, aggregated, and difficult to attribute to specific activities \citep{jessoeKnowledgeLessPower2014,tiefenbeckRealtimeFeedbackPromotes2019}. Energy bills often arrive only monthly or quarterly, weakening the link between everyday consumption choices and their monetary consequences \citep{chettySalienceTaxationTheory2009}. As a result, households may underweight the costs of electricity use and consume more than they would under fuller information \citep{alcocerSalienceBiasFramework2024}.

In this paper, we study whether digital information provision reduces household electricity consumption and whether such effects persist once initial engagement declines. To do so, we use a randomized field experiment with a staggered rollout in Upper Austria, the third most populous of Austria's nine federal provinces (about 1.5 million residents). In the experiment, households were randomly assigned either to an information condition, which received access to a smartphone application providing feedback on electricity use, costs, emissions, social comparisons, and gamified learning features, or to a control condition without access to the application. We link treatment assignment to high-frequency smart-meter data for up to 70 weeks after rollout and to detailed app-usage data from Google Analytics. The final analysis sample consists of 892 households, with 449 households in the information condition and 443 households in the control condition. Our main estimand is the intention-to-treat effect of being offered access to the application.


Our results show that access to the application reduces daily electricity consumption by around 4.5–4.7\%. This effect emerges within the first two weeks of rollout. Over the following year, point estimates remain consistently negative, but statistical significance is intermittent rather than continuous: the reduction is significant in most weeks through week 25 and only sporadically thereafter.

Furthermore, the reductions are concentrated in daytime consumption (07:00–22:00), the hours in which behavioural adjustments to electricity use are most feasible, while night-time consumption remains largely unchanged.

App usage declines sharply after the rollout, and households with higher observed engagement exhibit larger and more persistent reductions in electricity consumption. As engagement is measured after treatment assignment, these results describe engagement rather than causing it. Nevertheless, as both app usage and energy reductions decline over time, this correlation suggests that continuous exposure to information is necessary to maintain energy savings. 

A large empirical literature has examined the effectiveness of information-based and behavioral interventions in household energy consumption. Feedback interventions, social comparisons, and related tools have been shown to reduce electricity consumption in the short run, typically by 2-5\% on average, with some meta-analyses reporting effects of up to 7.4\% \citep{abrahamseReviewInterventionStudies2005,delmasInformationStrategiesEnergy2013,karlinEffectsFeedbackEnergy2015,andorBehavioralEconomicsEnergy2018,khannaMulticountryMetaanalysisRole2021,khannaBehavioralInformationMonetary2025}. At the same time, two questions remain open. First, much of the literature is based on relatively short intervention periods, making it difficult to assess whether treatment effects persist once the initial novelty of the intervention fades or user engagement declines. Several studies document attenuation over time, suggesting that at least part of the response may reflect temporary salience rather than durable behavioral adjustment \citep{allcottShortrunLongrunEffects2014}. Second, the behavioral mechanisms underlying these effects remain difficult to separate. Reductions in consumption may reflect persistent learning, temporary increases in attention and cost salience, or the formation of new habits \citep{lynhamWhyDoesRealtime2016,tiefenbeckOvercomingSalienceBias2018}.

Our design addresses these issues in two ways. Firstly, the extended observation period allows us to analyze the dynamics of treatment effects long after the initial rollout period. Secondly, app usage data allows us to establish whether changes in electricity consumption are linked to the level of engagement. The intervention comprises several components: feedback on consumption, costs and emissions; social comparison; and gamified learning features. Consequently, it is not possible to isolate the marginal causal effect of each component separately. While this is a limitation, it also reflects how digital energy interventions are commonly implemented in practice. Furthermore, recent meta-analytic evidence suggests that combinations of intervention components incorporating information, social and motivational elements are associated with stronger reductions in household energy consumption than the components alone \citep{khannaMulticountryMetaanalysisRole2021,khannaBehavioralInformationMonetary2025}. Therefore, our design is well suited to evaluating the effectiveness of a realistic digital intervention package that can be delivered at a low marginal cost once the necessary smart meter infrastructure is in place.


Our study contributes to the literature by providing evidence on the persistence of behavioral responses to digital information provision. We follow households for up to 70 weeks after rollout, substantially longer than most existing studies. This allows us to show that electricity consumption falls shortly after treatment and remains negative throughout, with statistical significance intermittent rather than continuous: detectable in most weeks through week 25 and only sporadically thereafter. This evidence speaks directly to the limited evidence on long-run effects in the existing literature. For example, \citet{khannaMulticountryMetaanalysisRole2021} report that the mean and median treatment durations in their sample were only 21.5 and 12 weeks, respectively.

The study also provides evidence on the time profile of treatment effects and user engagement. Existing studies show that feedback and information can reduce energy consumption, but effects vary substantially across settings, designs, and intervention features \citep{delmasInformationStrategiesEnergy2013,karlinEffectsFeedbackEnergy2015,andorBehavioralEconomicsEnergy2018,khannaBehavioralInformationMonetary2025}. A remaining question is whether reductions are sustained through learning, salience, habit formation, continued engagement, or some combination of these channels. 

By combining randomized treatment assignment with high-frequency consumption and app-usage data, we show that treatment effects are largest when engagement is highest and become less precisely estimated as app usage declines — descriptive evidence on the role of attention, though our design cannot fully separate the underlying mechanisms.

Finally, we evaluate a realistic digital intervention package that is relevant for policymakers and utilities. The application combines feedback on consumption, costs and emissions; social comparison; and gamified learning features. We therefore cannot isolate the marginal causal effect of each component. However, this bundled design reflects how digital energy interventions are commonly implemented in practice. Recent meta-analytic evidence also highlights the policy relevance of studying intervention packages, since real-world programs often combine informational, feedback-based, social, and motivational elements \citep{khannaBehavioralInformationMonetary2025}. Once smart-meter and app infrastructure is in place, such interventions can be delivered to additional households at very low marginal cost. Even modest household-level reductions can therefore generate meaningful aggregate effects when scaled across a large customer base.

The remainder of the paper is structured as follows. Section~\ref{sec:Conceptual_framework} outlines the conceptual framework. Section~\ref{sec:Experimental_design} describes the experimental design and data. Section~\ref{section:Results} presents the empirical results. Section~\ref{section:Conclusion} concludes.